\section{Behavioral constraints inside the search}\label{sec:behavior}

\subsection{The contract is not merely a final filter}

A complete design should support three interacting directions of information:
types constrain program shape; specifications propagate obligations backward
through that shape; observations and counterexamples constrain the remaining
choices. None of these replaces the final proof of the original contract.

Represent each behavioral obligation with its exact expression, local path
assumptions, program-hole dependencies, and proof status. A proof obligation
that mentions a completed candidate fragment may be attempted immediately. An
obligation dominated by an unconstrained function hole should usually wait for
more structure, or invoke a contract-decomposition rule rather than a generic
proof search that blindly guesses the function.

Proof tactics can assign metavariables. The controller must therefore choose
an explicit policy: freeze computational holes while checking a proposed
fragment, or allow a tactic to return computational assignments as a new
transactional synthesis proposal. An arbitrary proof worker must not silently
mutate unrelated program decisions. The proof and term engines share an
obligation graph, not ungoverned access to one mutable context.

\subsection{A bounded proof-service portfolio}

Begin with low-cost methods: conversion/reflexivity, existing local proofs,
constructor proofs, and a curated simplifier. Follow with indexed theorem
application, arithmetic or algebraic methods appropriate to the goal, and
bounded \code{grind} or Aesop invocations. Induction is a deliberate schema
choice, preferably coordinated with the program's recursion plan, rather than
a repeatedly retried final tactic.

Each call returns one of: a checked proof proposal; a checked refutation
proposal; a reduced residual goal with a proof of the reduction; unsupported;
or a resource/unknown diagnostic. A failed proof attempt does not falsify the
candidate. This preserves the sound distinction already present in Leant's
behavioral interface while allowing much richer search-directed use of the
contract~\cite{leant-behavior}.

Results must also preserve assumptions. A proof obtained under a temporary
branch guard can discharge only an obligation under that guard. A simplifier
that rewrites the target must return its equality or implication proof; its
pretty-printed residual expression alone is not enough to reconstruct the
original contract.

\subsection{Counterexample-guided synthesis}

For a universal contract
\[
 C(f)=\forall x:A,\ \mathsf{Pre}(x)\to\mathsf{Post}(x,f(x)),
\]
maintain a bank of typed inputs and their validity evidence. Candidate search
uses the bank as a necessary finite condition. A verifier tries to prove the
full contract, find a violating input, or report unknown. The loop is:
\begin{Code}
while resources remain:
  e := next candidate consistent with retained constraints
  if there is no such candidate:
      report exhaustion only for the declared explored grammar
  p := try proving the original C(e)
  if p is accepted:
      return Verified(e, p)
  outcome := seekCounterexample(e, C)
  if outcome is a replayed, valid counterexample x:
      add x; increment bank version; refine observation partitions
      propagate any certified necessary constraints
  else:
      retain e in an unresolved queue with its proof-search status
      continue with other candidates fairly
\end{Code}

A counterexample is not just a tuple of integers. It contains a value of the
actual input type, a proof of any dependent bounds and preconditions, and the
observed violation. For sound pruning, validate that violation in Lean or via
a separately certified evaluation procedure. A solver model containing an
out-of-bounds index does not refute a function defined only on \code{Fin n}.

If the user's request is only a finite conjunction of examples, a proof of
that conjunction completes the request. When a universal law was requested,
passing the entire current bank leaves a proof obligation. Exhaustive finite
domain verification can prove a universal statement only when the enumeration
has a checked coverage theorem and the relevant evaluation is sound.

\begin{proposition}[Progress of finite counterexample search]
For a finite candidate set of size $N$, suppose every rejected candidate has a
valid counterexample that is permanently retained, and candidate selection
never chooses a candidate violating the retained bank. Then at most $N$
distinct counterexample rejections occur before acceptance, verifier unknown,
or exhaustion of that candidate set.
\end{proposition}
\begin{proof}
Each new counterexample excludes the just-rejected candidate from all later
selection. Thus each rejection eliminates a previously selectable member of a
finite set. Unknown outcomes do not satisfy this argument and must not be
counted as counterexamples.
\end{proof}

The bound is not a practical runtime guarantee: constructing candidates and
proving or refuting them can dominate the work. Its purpose is to make the
candidate-retention and bank-update contracts explicit.

\subsection{Typed test generation}

Use a plugin registry of generators and quotation functions for supported
types. Constructors provide a starting grammar; indices and preconditions
provide filtering obligations. A generated \code{Fin n} value must include its
bound proof. A generated dependent pair first chooses the first field and
then generates a value of its instantiated second-field type.

For polymorphic functions, testing at several concrete types is useful
falsification evidence, not a proof of universal parametric behavior. Include
distinct concrete dictionary payloads, heterogeneous products, and asymmetric
values to avoid tests that accidentally identify meaningful alternatives.
Higher-order inputs can be sampled from finite tables or small lambda grammars
when their domains support this; such sampling remains incomplete.

Opaque, noncomputable, or effectful code needs an explicit evaluation policy.
A timeout, blocked reduction, or inability to quote a result is unknown. Never
convert it into false just to accelerate candidate rejection. The default
pure evaluator must not execute file, network, or process effects from an
unverified candidate.

\section{Certified abstraction and solver boundaries}\label{sec:abstraction}

\subsection{The exact condition for pruning a partial program}

Let $s$ be a typed sketch, $\Phi$ its valid path assumptions, and
$\Comp(s)$ the set of completions allowed by the current grammar. Let
$A(s,\Phi)$ be an abstract summary with concretization $\gamma$. A sufficient
soundness condition is
\begin{equation}\label{eq:abstract-soundness}
 \{\Obs(e,\rho):e\in\Comp(s),\ \rho\models\Phi\}
 \subseteq \gamma(A(s,\Phi)).
\end{equation}
The abstract summary overapproximates \emph{all} allowed completions and valid
inputs, not merely the completions sampled so far. For the point-observation
pruning rule below, also require an exhibited valid input $\rho_0\models\Phi$
and a defined observation for every admissible completion at that input.
Otherwise an empty input domain could make the universal contract vacuously
true while the observation set is empty. A plugin using whole-program summaries
instead must prove that every completion has a summary object.

Suppose $K$ is an abstractly expressible necessary condition for the original
contract. If a checked argument establishes
\[
 \gamma(A(s,\Phi))\cap K=\varnothing,
\]
then no completion of $s$ can satisfy that condition in the relevant setting.
The branch may be pruned with the certificate and its assumptions attached.
An underapproximation or a summary valid for just one chosen completion does
not support this conclusion.

\begin{proposition}[Sound sketch pruning]
Assume~\eqref{eq:abstract-soundness}, the valid-input and defined-observation
conditions above, and that every contract-satisfying completion has its
observation at $\rho_0$ in $K$. If the displayed intersection is empty,
then no completion represented by the sketch satisfies the contract under the
stated assumptions.
\end{proposition}
\begin{proof}
A satisfying completion has a defined observation at the exhibited input
$\rho_0$. That observation lies in $\gamma(A(s,\Phi))$ by soundness and in $K$
by necessity, contradicting the empty intersection.
\end{proof}

This elementary argument is often the missing step in an otherwise plausible
``use an SMT solver to prune'' design. Leant2 should encode it in the domain
plugin interface. A plugin lacking a certificate may still supply a ranking
hint; it cannot supply a semantic refutation.

\subsection{Domain plugins}

A useful plugin offers a typed reification procedure, an interpretation into
Lean, a relation between concrete and abstract states, transfer rules, and
proof builders for its claimed implications. Initial domains could cover
length and size, linear arithmetic, finite tags, bounded machine integers, and
simple multiset observations. The plugin may decline an unsupported provider
or return the top abstraction; refusal must not silently become a false fact.

The abstraction often needs assumptions about particular provider occurrences.
For example, a theorem describing a supplied callback's length behavior must
be instantiated with that exact callback and dictionary evidence. A theorem
for one occurrence is not transferable to a similarly printed occurrence from
a different branch. This preserves the useful candidate/evidence ownership
principles already present in the source repositories~\cite{leant-internals}.

The prototype experiment in Section~\ref{sec:experiments} gives a small complete
example: a typed fold-step AST has a provably exact affine length summary. The
Python search uses this summary, and Lean separately proves the summary rule
and the selected program's full behavior. The experiment uses no SMT solver.

\subsection{SMT as a proposal and proof-reconstruction service}

There are three different solver uses. A solver may propose choices for a
finite sketch, propose counterexample inputs, or prove a theory obligation
through a reconstructed certificate. Those uses have different validation
requirements; a single Boolean ``solver succeeded'' result is inadequate.

For synthesis, encode finite grammar choices, selected constants, and the
current examples. A satisfying assignment is decoded into a typed Lean plan
and checked. For counterexamples, encode the negation of the supported behavior
fragment and replay the decoded model against the exact candidate. For proofs,
reconstruct the result in Lean through a certificate checker or an established
proof-producing integration. Lean-SMT is a relevant optional integration point,
not a claim that all Lean propositions become SMT-decidable~\cite{lean-smt}.

Every theory translation has an explicit boundary. Naturals require
nonnegativity and the actual semantics of their operations. Bitvectors use
modular arithmetic. Arrays need their Lean bounds and extensionality semantics.
User-defined functions can be uninterpreted in an overapproximation, but models
of that overapproximation may not correspond to executable Lean values. A
proof under a sound weakening may still be useful; a countermodel needs replay.
Division, coercions, and finite-domain assumptions deserve dedicated regression
tests because superficially similar operations can have different edge cases.

An \code{unsat} answer for an underapproximate grammar excludes only that grammar.
An \code{unsat} answer for a sound overapproximation can support stronger
conclusions, but only after its encoding theorem and solver certificate are
accepted. Solver \code{unknown}, timeout, unsupported proof rules, and malformed
models remain distinct outcomes.

\subsection{Learning nogoods without leaking assumptions}

A failed branch can yield a useful nogood: these constructor choices, type
instantiations, and path assumptions cannot all occur in a successful
completion. Store the smallest justified dependency slice rather than the
entire failed syntax tree. Nevertheless, the nogood must identify the request,
contract, instance selections, relevant environment version, and grammar scope.
A solver unsat core is a proposed slice, not automatically a Lean theorem.

If failure means only that a proof method exhausted its budget, record a
heuristic penalty or a retry schedule. It is not a nogood. If failure follows
from a genuinely false finite observation, a checked lemma can sometimes reject
all candidates containing the same local decision under the same assumptions.
Generalization beyond that slice requires another proof.

\section{Equivalence, deduplication, and cache validity}\label{sec:equivalence}

\subsection{Observational agreement is provisional}

Let $X$ be the current finite example bank and define
\[
 e\sim_X e'\quad\Longleftrightarrow\quad
 \forall x\in X,\ \Obs(e,x)=\Obs(e',x).
\]
This relation groups candidates usefully. It is not universal extensional
equality, and in a higher-order setting it is not automatically a congruence
under every future context. Adding a counterexample can split an old class.

Store an observation class as a representative \emph{plus recoverable members}
and their generation recipes. When the bank changes, refine the partition or
resume deferred members. Permanent deletion is allowed only with an appropriate
semantic equivalence proof, a valid dominance argument for the requested
objective, or an explicitly incomplete fast-lane policy.

The finite experiment demonstrates the problem concretely. Identity and
reverse agree on the initial empty and singleton examples. Retaining only
identity deletes every successful reverse candidate in that initial class.
The next asymmetric list kills identity, leaving an apparently exhausted
search even though the original finite grammar contained a correct program.

\subsection{Definitional and proved equality}

Definitional equality is stronger evidence than agreement on tests, but its
use still depends on context and metavariable assignments. Compare completed
terms in a fixed environment or use a carefully abstracted template. Do not
use an equality test that mutates one branch's unknowns as a destructive global
deduplication key.

Proved equality can justify more aggressive normalization. A dependent e-graph
or rewrite system must track the type of each term, proofs of rewrites, and
transport through dependent parents. An ordinary untyped e-graph that merges
two children and leaves their parents unchanged can produce ill-typed terms.
Use homogeneous equality within a fixed type; when types change, construct
explicit transport or a verified heterogeneous-congruence step.

Proof irrelevance can help merge proof terms of the same proposition. It must
not merge operation dictionaries or arbitrary data records in \code{Type}.
Likewise, an equality between whole functions does not automatically justify a
cost dominance claim about their compiled implementations.

\subsection{Caching exact subproblems}

A positive cache entry should be a closed, parameterized template with an exact
type and its proof dependencies. Instantiation creates fresh metavariables and
rechecks applicability. Canonicalize locals by an explicit dependent telescope,
not by pretty-printed names. Include local definition values and the selected
instance terms where they affect meaning.

A negative cache entry is more demanding. It must record whether it is a
semantic refutation, exhaustive failure of a bounded grammar, or merely an
operational miss. A result obtained at cost bound $k$ cannot exclude a larger
bound. A failure with one set of providers cannot exclude another. An unknown
proof status should normally be cached only to avoid identical immediate
retries, with budget and method identity attached.

Hashing accelerates lookup; it does not confer authority. On a hit, compare or
validate the exact key material and replay the template. A request identity
includes the contract and policy, not just $T$. This is essential for the two
projection examples: they have the same function type and opposite behavioral
requirements.
